% ===================================================================
%  تابلوی شمارهٔ ۲ — بدن انسان. --- زنگ تفریح. بازی‌ها.
%
%  Traduction, faite en 2026, en persan d'aujourd'hui, sur l'ido de
%  texto/io. Regles de la colonne : voir 10-tablo-01.tex. Une seule
%  numerotation.
%
%  LE CORPS EST PERSAN, MAIS LES ORGANES SONT ARABES, ET LA FRONTIERE
%  EST D'UNE NETTETE QU'ON NE VOIT NULLE PART AILLEURS DANS LE
%  RELEVE. Tout ce qu'on peut montrer du doigt porte un nom iranien :
%  سر la tete, چهره le visage, پیشانی le front, گیجگاه la tempe, چشم
%  l'oeil, ابرو le sourcil, مژه le cil, بینی le nez, گونه la joue, گوش
%  l'oreille, دهان la bouche, لب la levre, دندان la dent, زبان la
%  langue, چانه le menton, گردن le cou, گلو la gorge, سینه la
%  poitrine, دنده la cote, پهلو le flanc, پشت le dos, شانه l'epaule,
%  شکم le ventre, بازو le bras, آرنج le coude, مچ le poignet, دست la
%  main, انگشت le doigt, ناخن l'ongle, ران la cuisse, زانو le genou,
%  ساق le mollet, قوزک la cheville, پا le pied, پاشنه le talon,
%  استخوان l'os, پوست la peau, خون le sang, مو le cheveu, مغز le
%  cerveau, شش le poumon, روده l'intestin, ماهیچه le muscle, رگ la
%  veine.
%
%  ET CE QU'IL FAUT OUVRIR POUR LE VOIR PORTE UN NOM ARABE : جمجمه le
%  crane, قلب le coeur, معده l'estomac, ساعد l'avant-bras. La raison
%  est une bibliotheque : la medecine du monde iranien s'est ECRITE en
%  arabe — Avicenne a compose le Qanun en arabe —, et le vocabulaire
%  de l'anatomie est revenu au persan dans la langue ou il avait ete
%  redige. Ce que l'enfant nomme est persan, ce que le medecin nomme
%  est arabe, et le partage suit exactement la peau.
%
%  LES CINQ SENS SONT CINQ NOMS FORMES SUR CINQ VERBES PERSANS PAR LE
%  MEME SUFFIXE : بینایی, شنوایی, بویایی, چشایی, بساوایی — de دیدن,
%  شنیدن, بوییدن, چشیدن, بساویدن. C'est exactement ce que le javanais
%  fait avec son prefixe pa- (pandeleng, pangrungu, pangambu,
%  pangecap, pangrasa) au meme alinea du meme tableau. Deux langues
%  sans parente, la meme elegance, et le francais et l'ido y opposent
%  cinq noms sans famille. Le tableau des sens est l'endroit du livret
%  ou une morphologie enseigne mieux que la lecon.
%
%  LA BICYCLETTE (37) SE DIT دوچرخه, « la deux-roues », ET C'EST LE
%  QUATRIEME MOT DU RELEVE POUR CETTE MACHINE. Le javanais dit pit, du
%  neerlandais fiets ; l'indonesien sepeda, du neerlandais velocipede ;
%  le francais bicyclette ; le persan ne prend rien et decrit. Quatre
%  langues, quatre solutions, une seule machine — et c'est la
%  bicyclette, arrivee partout dans les memes trente ans, qui mesure
%  le mieux ce que chaque langue fait d'une chose neuve.
%
%  ET LE PERSAN DECRIT SOUVENT LA OU IL POURRAIT EMPRUNTER : سرخرگ
%  l'artere, « la veine rouge » ; چوب‌پا les echasses (24), « le
%  bois-pied » ; گلخانه la serre (26), « la maison de fleurs » ;
%  فرفره la toupie (18) ; تیله les billes (21) ; بادبادک le cerf-volant
%  (33), « le petit vent-vent » ; قایم‌باشک le cache-cache (25). Mais
%  il emprunte le trapeze et le bilboquet, qui ne sont venus qu'avec
%  leur nom, et il DECRIT les quilles (16), qui ne sont pas venues du
%  tout. Trois conduites pour trois situations, et elles ne se
%  melangent pas.
% ===================================================================

%%K t02-apar-1 apar
\VUsaut{4.0mm}
\VUtitre{15.0pt}{60}{تابلوی شمارهٔ ۲}
\VUsaut{2.0mm}
\VUtitre{11.4pt}{0}{بدن انسان. --- زنگ تفریح.}
\VUtitre{11.4pt}{0}{بازی‌ها.}

%%K t02-tit-0 sub
\VUtitre{13.2pt}{0}{بدن انسان.}
\VUsaut{2.4mm}
\VUcentre{10.2pt}{0}{\textit{(درسی ساده از علوم طبیعی.)}}
\VUsaut{3.0mm}

%%K t02-tit-1 sub pos=t02-01-1
\VUpk{10.2pt}{40}{سر.}
\VUsaut{3.0mm}

%%K t02-01-1 p
1. --- بخش‌های اصلی بدن انسان اینهاست : \VUgras{سر}، \VUgras{تنه} و
\VUgras{اندام‌ها}.

%%K t02-02-1 p
2. --- سر دو بخش دارد : \VUgras{جمجمه} که \VUgras{مغز} را در بر گرفته
است و \VUgras{مو} آن را می‌پوشاند، و \VUgras{چهره}.

%%K t02-03-1 p
3. --- در نقاشی ما نه جمجمه دیده می‌شود و نه مغز ؛ اما بخش‌های مهم چهره
را بسیار روشن می‌بینیم.

%%K t02-04-1 p
4. --- نخست \VUgras{پیشانی} است که از اندیشهٔ نیرومند و همیشه در کار
نشان دارد. \VUgras{گیجگاه‌ها} بسیار آزادند. \VUgras{چشم} سرزنده و
گشوده است. \VUgras{ابروی} پرپشت و \VUgras{مژه‌های} بلند از آن نگهبانی
می‌کنند.

%%K t02-05-1 p
5. --- سپس \VUgras{بینی} و \VUgras{سوراخ‌های بینی} است. بینی برای
بوییدن و نفس کشیدن به کار ما می‌آید. در هر سوی بینی \VUgras{گونه‌ها}
هستند و دورتر \VUgras{گوش‌ها}. بخش پایین چهره از \VUgras{دهان} و
\VUgras{چانه} ساخته شده است.

%%K t02-06-1 p
6. --- آنها را چندان روشن نمی‌بینیم، زیرا \VUgras{سبیل} و \VUgras{خط
ریش} آنها را از ما پنهان می‌کنند. اما می‌دانیم که دهان دو \VUgras{لب}،
\VUgras{فک بالا} و \VUgras{فک پایین} با \VUgras{دندان‌هایشان}، و
\VUgras{زبان} دارد. با دهان سخن می‌گوییم و می‌خوریم. با زبان مزهٔ
خوراک را می‌چشیم.

%%K t02-07-1 p
7. --- چشم، گوش، بینی و دهان، مانند انگشتان، اندام‌های پنج حس‌اند :
\VUgras{بینایی}، \VUgras{شنوایی}، \VUgras{بویایی}، \VUgras{چشایی}،
\VUgras{بساوایی}.

%%K t02-08-1 p
8. --- پس سر یکی از جالب‌ترین بخش‌های بدن انسان است.

%%K t02-tit-2 sub
\VUsaut{4.4mm}
\VUpk{10.2pt}{40}{تنه.}
\VUsaut{3.0mm}

%%K t02-09-1 p
9. --- \VUgras{گردن} سر را به تنه می‌پیوندد. گردن \VUgras{گلو} و
\VUgras{پس‌گردن} را در بر دارد.

%%K t02-10-1 p
10. --- تنه اندام‌های بنیادی بسیاری هم دارد. در \VUgras{سینه}
\VUgras{شش‌ها} هستند که با آنها نفس می‌کشیم، و \VUgras{قلب} که کانون
زندگی است. هنگامی که قلب از تپیدن باز ایستد، موجود زنده می‌میرد.

%%K t02-11-1 p
11. --- \VUgras{دنده‌ها} سینه را نگه می‌دارند.

%%K t02-12-1 p
12. --- بخش‌های دیگر تنه \VUgras{پهلو}، \VUgras{پشت}، \VUgras{شانه‌ها}
و \VUgras{شکم} است. برای خوابیدن روی پهلو یا روی پشت می‌خوابیم، و
بارها را روی شانهٔ خود می‌بریم.

%%K t02-13-1 p
13. --- در شکم \VUgras{معده} و \VUgras{روده‌ها} هستند. آنها برای گواردن
خوراک به کار ما می‌آیند.

%%K t02-tit-3 sub
\VUsaut{4.4mm}
\VUpk{10.2pt}{40}{اندام‌ها.}
\VUsaut{3.0mm}

%%K t02-14-1 p
14. --- ما چهار \VUgras{اندام} داریم : دو \VUgras{بازو} و دو
\VUgras{پا}. با بازوها چیزها را برمی‌داریم، نگه می‌داریم، بلند می‌کنیم
و گرد می‌آوریم ؛ با پاها می‌ایستیم، راه می‌رویم و می‌دویم.

%%K t02-15-1 p
15. --- \VUgras{شانه} بازو را به بدن می‌پیوندد. \VUgras{ماهیچهٔ} بسیار
نیرومندی، یعنی دوسر، بازو را می‌جنباند، و \VUgras{آرنج} آن را به
\VUgras{ساعد} می‌پیوندد. \VUgras{مچ} بند \VUgras{دست} را می‌سازد، و
دست به \VUgras{انگشتان} و \VUgras{ناخن‌ها} پایان می‌یابد. روی دست
\VUgras{رگ} — و نیز سرخرگ — را باز می‌شناسیم که \VUgras{خون} در آن
روان است.

%%K t02-16-1 p
16. --- بخش‌های پا اینهاست : \VUgras{ران}، \VUgras{زانو} و \VUgras{پشت
زانو}، \VUgras{ساق} که \VUgras{گوشت} آن را \VUgras{پوست} پوشانده است،
\VUgras{قوزک} و \VUgras{کف پا}. \VUgras{استخوان‌های} پا بسیار کلفت و
بسیار نیرومندند. در انتهای پا \VUgras{انگشتان پا} هستند. بخش‌های دیگر
\VUgras{کف پا} و \VUgras{پاشنه} است.

%%K t02-17-1 p
17. --- چنین است شرح تقریباً کامل بدن انسان.

%%K t02-17-2 p
\textit{یادداشت.} --- \textit{آموزگار می‌تواند با کمک این عبارت :
« من در... (سر و مانند آن) درد دارم » همهٔ این واژه‌ها را از دیدگاه
خوبی یا بدی} \VUgras{تندرستی} \textit{دوباره به کار برد.}

%%K t02-tit-4 sub
\VUsaut{5.0mm}
\VUpk{10.2pt}{40}{ورزش.}
\VUsaut{2.4mm}
\VUcentre{10.2pt}{0}{\textit{(روایت یکی از دانش‌آموزان.)}}
\VUsaut{3.0mm}

%%K t02-18-1 p
18. --- برای آنکه نرم، چابک و تندرست باشیم، باید پاها و بازوهای خود را
ورزش دهیم. از این رو بازی می‌کنیم و \VUgras{ورزش} می‌کنیم.

%%K t02-19-1 p
19. --- در طول هفته، این دو تمرین در حیاط دبیرستان انجام می‌شود (تابلوی
شمارهٔ ۱ را ببینید). اما پنجشنبه و یکشنبه به چمنزار بزرگی می‌رویم، جایی
که برای دویدن و خوش گذراندن جا داریم.

%%K t02-20-1 p
20. --- آموزگاران و پدر و مادرهای ما گاهی ما را تا آنجا همراهی می‌کنند ؛
اما \VUgras{آموزگار ورزش}\textsuperscript{(12)} است که تمرین‌ها را
راهبری می‌کند.

%%K t02-21-1 p
21. --- در چمنزار، سمت چپ درِ ورودی، \VUgras{وسایل
ورزش}\textsuperscript{(1)} هست ؛ ما از \VUgras{نردبان}\textsuperscript{(2)}
بالا می‌رویم، روی \VUgras{حلقه‌ها}\textsuperscript{(3)} و روی
\VUgras{ترپز}\textsuperscript{(4)} وارونه می‌شویم، و خود را با دست تا
بالای \VUgras{نردبان طنابی}\textsuperscript{(5)} یا \VUgras{طناب
گره‌دار}\textsuperscript{(6)} می‌کشیم.

%%K t02-22-1 p
22. --- در همان هنگام، رفیقان ما از روی \VUgras{اسب
چوبی}\textsuperscript{(7)} یا \VUgras{میله‌های
موازی}\textsuperscript{(11)} می‌پرند. دیگران \VUgras{مشت‌زنی}\textsuperscript{(8)}
یا \VUgras{شمشیربازی}\textsuperscript{(9)} می‌کنند، با ماسک و
\VUgras{فلوره.} سرانجام، دیگران \VUgras{وزنه‌برداری}\textsuperscript{(10)}
می‌کنند.

%%K t02-tit-5 sub
\VUsaut{4.6mm}
\VUpk{10.2pt}{40}{بازی‌ها.}
\VUsaut{3.0mm}

%%K t02-23-1 p
23. --- گرداگرد ورزشکاران، پسرها و دخترها \VUgras{بازی‌های} گوناگون
می‌کنند. دخترها با \VUgras{حلقهٔ بزرگ}\textsuperscript{(14)} خود خوش
می‌گذرانند ؛ با \VUgras{طناب}\textsuperscript{(13)} می‌پرند یا برادران و
پسرعموهای خود را به بازی \VUgras{کروکه}\textsuperscript{(15)} می‌خوانند.
شادمان می‌شوند که \VUgras{گوی} حریف را با \VUgras{چکش چوبی} خود دور
بزنند، درست هنگامی که او می‌پندارد به آسانی از طاق کوچک خواهد گذشت!

%%K t02-24-1 p
24. --- پسرها بازی‌های ماهیچه‌ای‌تر را بیشتر دوست دارند. با شوق
\VUgras{میله‌های بازی}\textsuperscript{(16)} را می‌زنند و با
\VUgras{گوی}\textsuperscript{(17)} چابکانه می‌اندازند. بازی با
\VUgras{فرفره}\textsuperscript{(18)} و \VUgras{بیلبوکه}\textsuperscript{(19)}
یا حتی \VUgras{بازی قورباغه}\textsuperscript{(20)} را هم خوار
نمی‌شمارند. اما بازی‌های دلخواه آنان \VUgras{تیله}\textsuperscript{(21)}
و \VUgras{توپ دیواری}\textsuperscript{(22)} است، زیرا دیوار
\VUgras{گلخانه}\textsuperscript{(26)} برای برگرداندن توپ سبک بسیار
مناسب است.

%%K t02-25-1 p
25. --- یوهانس کوچک که با \VUgras{چوب‌پا}\textsuperscript{(24)} راه
می‌رود، بسیار چابک است و خوب می‌داند چگونه تعادل خود را نگه دارد. نزدیک
او، چند رفیق در آن گلخانه و پشت \VUgras{پرچین} \VUgras{قایم‌باشک}\textsuperscript{(25)}
بازی می‌کنند. دیگران \VUgras{بازی حدس ضربه}\textsuperscript{(28)} یا
\VUgras{بدمینتون}\textsuperscript{(29)} را بیشتر می‌پسندند، که توپش از
یک \VUgras{راکت} به راکت دیگر می‌پرد.

%%K t02-noto-1 noto
\VUnotes{143.2mm}{%
(*) بازی‌های کودکان بارها نام‌هایی سراسر ملی دارند. ما کوشیدیم آنها را
تا آنجا که می‌شد در ایدو نام بگذاریم، گاه با الهام از خود کاری که
بازی را نشان می‌دهد، گاه با برگزیدن نام ملی این کشور یا آن کشور، هرگاه
آن نام چندان نادرست نباشد. برای نمونه : \VUgras{barelo-ludo} (آلمانی و
فرانسوی) همان \VUgras{rano-ludo} \textsuperscript{(20)} (انگلیسی) است،
که در آن بازیکنان می‌کوشند قرص‌های گرد خود را از میان دهان قورباغهٔ
فلزی بگذرانند که دهان‌گشوده روی جعبهٔ سوراخ‌دار (بشکه) ایستاده است.
\VUgras{shultro-salto}\textsuperscript{(30)} تنها در این نکته با
\VUgras{dorso-salto} فرق دارد : برای پریدن از روی سر، پرندگان دست‌های
خود را بر شانهٔ همبازی می‌گذارند، حال آنکه در « \VUgras{dorso-salto} »
تنها بر پشت او می‌گذارند. --- یا نیز چند تن از شرکت‌کنندگان خم می‌شوند
و دیگران از روی پشت آنان می‌پرند و دست بر شانه می‌گذارند ؛ گروه نخست
باید ضربه را بی‌آنکه فرو بنشیند تاب آورد، و گروه دوم باید بی‌آنکه
تعادل خود را از دست بدهد بپرد (خود تابلو را ببینید).%
}

%%K t02-26-1 p
26. --- در میان چمنزار دو درخت هست. زیر یکی از آن دو، هفت یا هشت پسر
بازی \VUgras{پرش از روی شانه}\textsuperscript{(30)} (*) را برپا
کرده‌اند. این بازی در همهٔ کشورها شناخته نیست ؛ اما همه می‌دانند
\VUgras{پرش از روی پشت} چیست، که این بار بازی نمی‌کنند.

%%K t02-tit-6 sub
\VUsaut{5.0mm}
\VUpk{10.2pt}{40}{بازی‌ها در هوای آزاد.}
\VUsaut{3.4mm}

%%K t02-27-1 p
27. \VUgras{بازی چشم‌بسته}\textsuperscript{(31)} نیز بسیار سرگرم‌کننده
است ؛ دخترها و پسرها به نوبت \VUgras{چشم‌بند} را بر چشم خود می‌بندند و
ناشیانی را می‌گیرند که می‌گذارند گرفته شوند.

%%K t02-28-1 p
28. --- در میان چمنزار، این هم بازی بسیار فرانسوی اما اندکی خشن :
\VUgras{بازی خرس}\textsuperscript{(32)}، که از بازی آرام
\VUgras{بادبادک}\textsuperscript{(33)} و حتی از بازی انگلیسی
\VUgras{تنیس}\textsuperscript{(34)} پرهیاهوتر است.

%%K t02-29-1 p
29. --- این ورزش آخری شادی دانش‌آموزان بزرگ است. آموزگاران ما نیز با
شوق آن را انجام می‌دهند. این ورزش چالاکی و مهارت بسیار می‌خواهد و مرا
بسیار خوش می‌آید. بازی \VUgras{تاب}\textsuperscript{(35)} را به دخترها
وامی‌گذارم.

%%K t02-30-1 p
30. --- بازی انگلیسی دیگری را دوست ندارم که شاید خطرناک شود.

%%K t02-30-1b p
منظورم \VUgras{فوتبال}\textsuperscript{(36)} است که برادر بزرگ‌ترم را
شاد می‌کند. \VUgras{بازی خط}\textsuperscript{(38)} کهن خودمان را بیشتر
می‌پسندم که همان اندازه بهداشتی و به راستی کم‌خشونت‌تر است.

%%K t02-tit-7 sub
\VUsaut{4.6mm}
\VUpk{10.2pt}{40}{دوچرخه‌سواری و بادکنک.}
\VUsaut{3.0mm}

%%K t02-31-1 p
31. --- با شوق نیز چمنزار را با \VUgras{دوچرخه}\textsuperscript{(37)}
دور می‌زنم.

%%K t02-32-1 p
32. --- سال آینده \VUgras{اسب‌سواری}\textsuperscript{(39)} خواهم آموخت.
چه زیباست اسبی که با ناز گام برمی‌دارد یا یورتمه می‌رود، و چهارنعل از
روی مانع‌ها می‌پرد!

%%K t02-33-1 p
33. --- در تابستان \VUgras{شنا}\textsuperscript{(40)} می‌آموزیم. با
لذت در آب روشن و خنک رودخانه شیرجه می‌رویم و آب‌تنی می‌کنیم. گاهی در
پایان \VUgras{بادکنک}\textsuperscript{(41)} کاغذی‌ای رها می‌کنیم که به
شکوه در هوا بالا می‌رود، همین که از هوای گرم پر شود.

%%K t02-34-1 p
34. --- بازی‌های بسیار دیگری هم هست، اما شمردن همهٔ آنها را به پایان
نمی‌رسانم ؛ و از این خلاصه دانسته می‌شود که در پنجشنبه‌ها در چمنزار
زیبای ما کسی دلتنگ نمی‌شود.

%%K t02-34-2 p
یادداشت. --- \textit{آموزگار به آسانی این فصل را کامل خواهد کرد ؛ با
شرح دقیق‌تر هر یک از مهم‌ترین بازی‌ها.}
